
%%%%%%%%%%%%%%%%%%%%%%%%%%%%%%%%%%%%%%%%%%%%%%%%%%%%%%%%%%%%%%%%%%%%%%%%%%%%%%

%\section{Some Auxiliary Lemmas}
%\label{sec: some auxiliary lemmas}
%
%%%%%%%%%%%%%%%%%%%%%%%%%%%%%%%%%%%%%%%%%%%%%%%%%%%%%%%%%%%%%%%%%%%%%%%%%%%%%%

%\section{Some Auxiliary Lemmas}
%\label{sec: some auxiliary lemmas}
%
%%%%%%%%%%%%%%%%%%%%%%%%%%%%%%%%%%%%%%%%%%%%%%%%%%%%%%%%%%%%%%%%%%%%%%%%%%%%%%

%\section{Some Auxiliary Lemmas}
%\label{sec: some auxiliary lemmas}
%
%%%%%%%%%%%%%%%%%%%%%%%%%%%%%%%%%%%%%%%%%%%%%%%%%%%%%%%%%%%%%%%%%%%%%%%%%%%%%%

%\section{Some Auxiliary Lemmas}
%\label{sec: some auxiliary lemmas}
%\input{4_auxiliary_lemmas.tex}

%%%%%%%%



In this section we show that, for Condition~{\MixCond},
we only need to consider $b$'s that are odd-prime factors of $n$
and satisfy $b\leq c_{max}$,
where $c_{max}$ denotes the maximum absolute value of concentrations in $C$.
These properties play an important role in the sufficiency proof of Condition~{\MixCond} in 
Theorem~\ref{thm: miscibility characterization}(a).
Additionally, they lead to an efficient algorithm for testing perfect mixability
(see Section~\ref{sec: polynomial running time}), 
thus showing the first part of Theorem~\ref{thm: miscibility characterization}(c). 

%%%%%%%%

\begin{lemma}\label{lem: uniform wrt coprime numbers}
Let $b,c\in\posintegers$ and $A\ingroundset \integers$.
(a) If $A$ is $bc$-congruent then $A$ is also $b$-congruent.
(b) If $\mygcd(b,c) = 1$ and $A$ is both $b$-congruent and $c$-congruent
	then $A$ is $bc$-congruent.
\end{lemma}

\begin{proof}
Part~(a) is trivial, because $x\equiv y\pmod{bc}$ implies that
$x\equiv y\pmod{b}$.
Part~(b) is also simple: Suppose that $x\equiv y\pmod{b}$
and $x\equiv y\pmod{c}$. This means that $b|(x-y)$ and $c|(x-y)$.
This, since $b,c$ are co-prime, implies that $(bc)|(x-y)$, 
which is equivalent to $x\equiv y\pmod{bc}$.
\end{proof}

%%%%%%%%%

\begin{lemma}\label{lem: mc powers of primes}
If Condition~{\MixCond} holds for all $b\in\posintegers$ that are a power of an odd prime
then it holds for all odd $b\in\posintegers$.
\end{lemma}

\begin{proof}
Assume that condition~{\MixCond} holds for all $b$ that are odd prime powers. 
Let $b'\in\posintegers$ be odd, with factorization
$b' = p_1^{\gamma_1} ... p_k^{\gamma_k}$,
for different odd primes $p_1 , ...,p_k$, and suppose that $C$ is $b'$-congruent.
Then, by Lemma~\ref{lem: uniform wrt coprime numbers}(a),
$C$ is also  $p_i^{\gamma_i}$-congruent for all $i$.
Since condition~{\MixCond} holds for $p_i^{\gamma_i}$,
this implies that $C\cup\braced{\mu}$ is $p_i^{\gamma_i}$-congruent for all $i$.
By repeated application of Lemma~\ref{lem: uniform wrt coprime numbers}(b), we then obtain that
$C\cup\braced{\mu}$ is $b'$-congruent as well. So Condition~{\MixCond} holds for $b'$.
\end{proof}

%%%%%%%%%

\begin{lemma}\label{lem: b's co-prime to n}
Let $b\in\posintegers$ be odd. If $\gcd(b,n) = 1$ then
Condition~{\MixCond} holds for $b$.
\end{lemma}

\begin{proof}
Assume that $C$ is $b$-congruent. 
By Observation~\ref{obs: offsetting does not affect reachability}, without loss of generality
we can assume that all numbers in $C$ are multiples of $b$. 
(Otherwise we can consider $C' = C-c$, for an arbitrary $c\in C$, 
because Condition~{\MixCond} is not affected by offsetting $C$.)
Thus $\mysum(C) = b\beta$, for some $\beta\in\integers$,
which gives us that $\mu = \mysum(C)/n = b\beta/n$.
As $\mu$ is integer and $\gcd(b,n) = 1$, $\beta$ must be a multiple of $n$.
We can thus conclude that $\mu$ is a multiple of $b$.
This means that $C\cup \braced{\mu}$ is $b$-congruent, proving that
Condition~{\MixCond} holds for $b$.
\end{proof}

%%%%%%%%%

\begin{corollary}\label{cor: mc prime power factors}
	Let $c_{max}$ be the maximum absolute value of concentrations in $C$. 
	If Condition~{\MixCond} holds for all $b\in\posintegers$ with $b\leq c_{max}$
	that are powers of odd prime factors of $n$,
	then Condition~{\MixCond} holds for all odd $b$.
\end{corollary}

%%%%%%%%%


To substantiate Corollary~\ref{cor: mc prime power factors},
note that $C$ satisfies Condition~$\MixCond$ 
for all odd $b\in\posintegers$ larger than $c_{max}$.
This holds because for each such $b$, $c\mod b=c$ for all $c\in C$, 
so the remainders of the (different) concentrations in $C$ modulo $b$ are all different. 
(In the trivial case where 
all concentrations in $C$ are equal, $\mu$ is also equal and thus
$C$ satisfies Condition~{\MixCond}; such $C$ is actually perfectly mixed already.)



%%%%%%%%



In this section we show that, for Condition~{\MixCond},
we only need to consider $b$'s that are odd-prime factors of $n$
and satisfy $b\leq c_{max}$,
where $c_{max}$ denotes the maximum absolute value of concentrations in $C$.
These properties play an important role in the sufficiency proof of Condition~{\MixCond} in 
Theorem~\ref{thm: miscibility characterization}(a).
Additionally, they lead to an efficient algorithm for testing perfect mixability
(see Section~\ref{sec: polynomial running time}), 
thus showing the first part of Theorem~\ref{thm: miscibility characterization}(c). 

%%%%%%%%

\begin{lemma}\label{lem: uniform wrt coprime numbers}
Let $b,c\in\posintegers$ and $A\ingroundset \integers$.
(a) If $A$ is $bc$-congruent then $A$ is also $b$-congruent.
(b) If $\mygcd(b,c) = 1$ and $A$ is both $b$-congruent and $c$-congruent
	then $A$ is $bc$-congruent.
\end{lemma}

\begin{proof}
Part~(a) is trivial, because $x\equiv y\pmod{bc}$ implies that
$x\equiv y\pmod{b}$.
Part~(b) is also simple: Suppose that $x\equiv y\pmod{b}$
and $x\equiv y\pmod{c}$. This means that $b|(x-y)$ and $c|(x-y)$.
This, since $b,c$ are co-prime, implies that $(bc)|(x-y)$, 
which is equivalent to $x\equiv y\pmod{bc}$.
\end{proof}

%%%%%%%%%

\begin{lemma}\label{lem: mc powers of primes}
If Condition~{\MixCond} holds for all $b\in\posintegers$ that are a power of an odd prime
then it holds for all odd $b\in\posintegers$.
\end{lemma}

\begin{proof}
Assume that condition~{\MixCond} holds for all $b$ that are odd prime powers. 
Let $b'\in\posintegers$ be odd, with factorization
$b' = p_1^{\gamma_1} ... p_k^{\gamma_k}$,
for different odd primes $p_1 , ...,p_k$, and suppose that $C$ is $b'$-congruent.
Then, by Lemma~\ref{lem: uniform wrt coprime numbers}(a),
$C$ is also  $p_i^{\gamma_i}$-congruent for all $i$.
Since condition~{\MixCond} holds for $p_i^{\gamma_i}$,
this implies that $C\cup\braced{\mu}$ is $p_i^{\gamma_i}$-congruent for all $i$.
By repeated application of Lemma~\ref{lem: uniform wrt coprime numbers}(b), we then obtain that
$C\cup\braced{\mu}$ is $b'$-congruent as well. So Condition~{\MixCond} holds for $b'$.
\end{proof}

%%%%%%%%%

\begin{lemma}\label{lem: b's co-prime to n}
Let $b\in\posintegers$ be odd. If $\gcd(b,n) = 1$ then
Condition~{\MixCond} holds for $b$.
\end{lemma}

\begin{proof}
Assume that $C$ is $b$-congruent. 
By Observation~\ref{obs: offsetting does not affect reachability}, without loss of generality
we can assume that all numbers in $C$ are multiples of $b$. 
(Otherwise we can consider $C' = C-c$, for an arbitrary $c\in C$, 
because Condition~{\MixCond} is not affected by offsetting $C$.)
Thus $\mysum(C) = b\beta$, for some $\beta\in\integers$,
which gives us that $\mu = \mysum(C)/n = b\beta/n$.
As $\mu$ is integer and $\gcd(b,n) = 1$, $\beta$ must be a multiple of $n$.
We can thus conclude that $\mu$ is a multiple of $b$.
This means that $C\cup \braced{\mu}$ is $b$-congruent, proving that
Condition~{\MixCond} holds for $b$.
\end{proof}

%%%%%%%%%

\begin{corollary}\label{cor: mc prime power factors}
	Let $c_{max}$ be the maximum absolute value of concentrations in $C$. 
	If Condition~{\MixCond} holds for all $b\in\posintegers$ with $b\leq c_{max}$
	that are powers of odd prime factors of $n$,
	then Condition~{\MixCond} holds for all odd $b$.
\end{corollary}

%%%%%%%%%


To substantiate Corollary~\ref{cor: mc prime power factors},
note that $C$ satisfies Condition~$\MixCond$ 
for all odd $b\in\posintegers$ larger than $c_{max}$.
This holds because for each such $b$, $c\mod b=c$ for all $c\in C$, 
so the remainders of the (different) concentrations in $C$ modulo $b$ are all different. 
(In the trivial case where 
all concentrations in $C$ are equal, $\mu$ is also equal and thus
$C$ satisfies Condition~{\MixCond}; such $C$ is actually perfectly mixed already.)



%%%%%%%%



In this section we show that, for Condition~{\MixCond},
we only need to consider $b$'s that are odd-prime factors of $n$
and satisfy $b\leq c_{max}$,
where $c_{max}$ denotes the maximum absolute value of concentrations in $C$.
These properties play an important role in the sufficiency proof of Condition~{\MixCond} in 
Theorem~\ref{thm: miscibility characterization}(a).
Additionally, they lead to an efficient algorithm for testing perfect mixability
(see Section~\ref{sec: polynomial running time}), 
thus showing the first part of Theorem~\ref{thm: miscibility characterization}(c). 

%%%%%%%%

\begin{lemma}\label{lem: uniform wrt coprime numbers}
Let $b,c\in\posintegers$ and $A\ingroundset \integers$.
(a) If $A$ is $bc$-congruent then $A$ is also $b$-congruent.
(b) If $\mygcd(b,c) = 1$ and $A$ is both $b$-congruent and $c$-congruent
	then $A$ is $bc$-congruent.
\end{lemma}

\begin{proof}
Part~(a) is trivial, because $x\equiv y\pmod{bc}$ implies that
$x\equiv y\pmod{b}$.
Part~(b) is also simple: Suppose that $x\equiv y\pmod{b}$
and $x\equiv y\pmod{c}$. This means that $b|(x-y)$ and $c|(x-y)$.
This, since $b,c$ are co-prime, implies that $(bc)|(x-y)$, 
which is equivalent to $x\equiv y\pmod{bc}$.
\end{proof}

%%%%%%%%%

\begin{lemma}\label{lem: mc powers of primes}
If Condition~{\MixCond} holds for all $b\in\posintegers$ that are a power of an odd prime
then it holds for all odd $b\in\posintegers$.
\end{lemma}

\begin{proof}
Assume that condition~{\MixCond} holds for all $b$ that are odd prime powers. 
Let $b'\in\posintegers$ be odd, with factorization
$b' = p_1^{\gamma_1} ... p_k^{\gamma_k}$,
for different odd primes $p_1 , ...,p_k$, and suppose that $C$ is $b'$-congruent.
Then, by Lemma~\ref{lem: uniform wrt coprime numbers}(a),
$C$ is also  $p_i^{\gamma_i}$-congruent for all $i$.
Since condition~{\MixCond} holds for $p_i^{\gamma_i}$,
this implies that $C\cup\braced{\mu}$ is $p_i^{\gamma_i}$-congruent for all $i$.
By repeated application of Lemma~\ref{lem: uniform wrt coprime numbers}(b), we then obtain that
$C\cup\braced{\mu}$ is $b'$-congruent as well. So Condition~{\MixCond} holds for $b'$.
\end{proof}

%%%%%%%%%

\begin{lemma}\label{lem: b's co-prime to n}
Let $b\in\posintegers$ be odd. If $\gcd(b,n) = 1$ then
Condition~{\MixCond} holds for $b$.
\end{lemma}

\begin{proof}
Assume that $C$ is $b$-congruent. 
By Observation~\ref{obs: offsetting does not affect reachability}, without loss of generality
we can assume that all numbers in $C$ are multiples of $b$. 
(Otherwise we can consider $C' = C-c$, for an arbitrary $c\in C$, 
because Condition~{\MixCond} is not affected by offsetting $C$.)
Thus $\mysum(C) = b\beta$, for some $\beta\in\integers$,
which gives us that $\mu = \mysum(C)/n = b\beta/n$.
As $\mu$ is integer and $\gcd(b,n) = 1$, $\beta$ must be a multiple of $n$.
We can thus conclude that $\mu$ is a multiple of $b$.
This means that $C\cup \braced{\mu}$ is $b$-congruent, proving that
Condition~{\MixCond} holds for $b$.
\end{proof}

%%%%%%%%%

\begin{corollary}\label{cor: mc prime power factors}
	Let $c_{max}$ be the maximum absolute value of concentrations in $C$. 
	If Condition~{\MixCond} holds for all $b\in\posintegers$ with $b\leq c_{max}$
	that are powers of odd prime factors of $n$,
	then Condition~{\MixCond} holds for all odd $b$.
\end{corollary}

%%%%%%%%%


To substantiate Corollary~\ref{cor: mc prime power factors},
note that $C$ satisfies Condition~$\MixCond$ 
for all odd $b\in\posintegers$ larger than $c_{max}$.
This holds because for each such $b$, $c\mod b=c$ for all $c\in C$, 
so the remainders of the (different) concentrations in $C$ modulo $b$ are all different. 
(In the trivial case where 
all concentrations in $C$ are equal, $\mu$ is also equal and thus
$C$ satisfies Condition~{\MixCond}; such $C$ is actually perfectly mixed already.)



%%%%%%%%



In this section we show that, for Condition~{\MixCond},
we only need to consider $b$'s that are odd-prime factors of $n$
and satisfy $b\leq c_{max}$,
where $c_{max}$ denotes the maximum absolute value of concentrations in $C$.
These properties play an important role in the sufficiency proof of Condition~{\MixCond} in 
Theorem~\ref{thm: miscibility characterization}(a).
Additionally, they lead to an efficient algorithm for testing perfect mixability
(see Section~\ref{sec: polynomial running time}), 
thus showing the first part of Theorem~\ref{thm: miscibility characterization}(c). 

%%%%%%%%

\begin{lemma}\label{lem: uniform wrt coprime numbers}
Let $b,c\in\posintegers$ and $A\ingroundset \integers$.
(a) If $A$ is $bc$-congruent then $A$ is also $b$-congruent.
(b) If $\mygcd(b,c) = 1$ and $A$ is both $b$-congruent and $c$-congruent
	then $A$ is $bc$-congruent.
\end{lemma}

\begin{proof}
Part~(a) is trivial, because $x\equiv y\pmod{bc}$ implies that
$x\equiv y\pmod{b}$.
Part~(b) is also simple: Suppose that $x\equiv y\pmod{b}$
and $x\equiv y\pmod{c}$. This means that $b|(x-y)$ and $c|(x-y)$.
This, since $b,c$ are co-prime, implies that $(bc)|(x-y)$, 
which is equivalent to $x\equiv y\pmod{bc}$.
\end{proof}

%%%%%%%%%

\begin{lemma}\label{lem: mc powers of primes}
If Condition~{\MixCond} holds for all $b\in\posintegers$ that are a power of an odd prime
then it holds for all odd $b\in\posintegers$.
\end{lemma}

\begin{proof}
Assume that condition~{\MixCond} holds for all $b$ that are odd prime powers. 
Let $b'\in\posintegers$ be odd, with factorization
$b' = p_1^{\gamma_1} ... p_k^{\gamma_k}$,
for different odd primes $p_1 , ...,p_k$, and suppose that $C$ is $b'$-congruent.
Then, by Lemma~\ref{lem: uniform wrt coprime numbers}(a),
$C$ is also  $p_i^{\gamma_i}$-congruent for all $i$.
Since condition~{\MixCond} holds for $p_i^{\gamma_i}$,
this implies that $C\cup\braced{\mu}$ is $p_i^{\gamma_i}$-congruent for all $i$.
By repeated application of Lemma~\ref{lem: uniform wrt coprime numbers}(b), we then obtain that
$C\cup\braced{\mu}$ is $b'$-congruent as well. So Condition~{\MixCond} holds for $b'$.
\end{proof}

%%%%%%%%%

\begin{lemma}\label{lem: b's co-prime to n}
Let $b\in\posintegers$ be odd. If $\gcd(b,n) = 1$ then
Condition~{\MixCond} holds for $b$.
\end{lemma}

\begin{proof}
Assume that $C$ is $b$-congruent. 
By Observation~\ref{obs: offsetting does not affect reachability}, without loss of generality
we can assume that all numbers in $C$ are multiples of $b$. 
(Otherwise we can consider $C' = C-c$, for an arbitrary $c\in C$, 
because Condition~{\MixCond} is not affected by offsetting $C$.)
Thus $\mysum(C) = b\beta$, for some $\beta\in\integers$,
which gives us that $\mu = \mysum(C)/n = b\beta/n$.
As $\mu$ is integer and $\gcd(b,n) = 1$, $\beta$ must be a multiple of $n$.
We can thus conclude that $\mu$ is a multiple of $b$.
This means that $C\cup \braced{\mu}$ is $b$-congruent, proving that
Condition~{\MixCond} holds for $b$.
\end{proof}

%%%%%%%%%

\begin{corollary}\label{cor: mc prime power factors}
	Let $c_{max}$ be the maximum absolute value of concentrations in $C$. 
	If Condition~{\MixCond} holds for all $b\in\posintegers$ with $b\leq c_{max}$
	that are powers of odd prime factors of $n$,
	then Condition~{\MixCond} holds for all odd $b$.
\end{corollary}

%%%%%%%%%


To substantiate Corollary~\ref{cor: mc prime power factors},
note that $C$ satisfies Condition~$\MixCond$ 
for all odd $b\in\posintegers$ larger than $c_{max}$.
This holds because for each such $b$, $c\mod b=c$ for all $c\in C$, 
so the remainders of the (different) concentrations in $C$ modulo $b$ are all different. 
(In the trivial case where 
all concentrations in $C$ are equal, $\mu$ is also equal and thus
$C$ satisfies Condition~{\MixCond}; such $C$ is actually perfectly mixed already.)

